% !TEX program = xelatex
\documentclass[12pt,a4paper,openany,fontset=fandol]{ctexrep}
\usepackage{amsmath,amssymb,mathtools,physics,bm}
\usepackage{geometry,graphicx,booktabs,array,longtable}
\usepackage{xcolor,hyperref}
\setCJKmainfont[AutoFakeBold=2]{gbsn00lp.ttf}
\setCJKsansfont[AutoFakeBold=2]{gbsn00lp.ttf}
\setCJKmonofont[AutoFakeBold=2]{gkai00mp.ttf}
\geometry{margin=2.5cm}
\hypersetup{colorlinks=true,linkcolor=blue!55!black,urlcolor=blue!55!black,
  pdftitle={量子力学中的常见效应：物理图像、数学推导与应用}}
\setlength{\parindent}{2em}
\setlength{\parskip}{0.25em}
\allowdisplaybreaks[2]
\newcommand{\ii}{\mathrm{i}}
\newcommand{\ee}{\mathrm{e}}
\newcommand{\kb}{k_{\mathrm B}}
\newcommand{\Physical}{\subsection*{物理图像}}
\newcommand{\Model}{\subsection*{模型与哈密顿量}}
\newcommand{\Derivation}{\subsection*{关键推导}}
\newcommand{\Conclusion}{\subsection*{核心结论}}
\newcommand{\Application}{\subsection*{应用与联系}}
\newcommand{\trans}[1]{\par\medskip\noindent\textcolor{blue!55!black}{\textbf{过渡：}}#1\par\medskip}

\title{\Huge 量子力学中的常见效应\\[0.6em]
  \Large 物理图像、数学推导与应用}
\author{系统学习讲义}
\date{\today}

\begin{document}
\maketitle
\begin{abstract}
本讲义以少数共同结构串联量子力学中的典型效应：边界条件与量子化、算符不对易与涨落、
振幅叠加与干涉、交换对称性与量子统计、外场耦合与能级劈裂、规范势与几何相位、
多体关联以及拓扑不变量。推导以本科量子力学为起点，对量子光学、凝聚态与量子场论部分
给出标准模型下的核心论证。符号约定为：$\hbar=h/(2\pi)$，电子电荷为$-e$（$e>0$）。
\end{abstract}
\tableofcontents

\chapter{量子效应的共同来源}
所谓“量子效应”并非一组互不相关的奇观，而是同一套数学原则在不同体系中的表现。

\section{态、振幅与测量}
纯态是希尔伯特空间中的射线$|\psi\rangle$；在位置表象中
$\psi(\bm r)=\langle\bm r|\psi\rangle$。线性方程允许
$|\psi\rangle=c_1|\psi_1\rangle+c_2|\psi_2\rangle$，而 Born 规则把概率写成振幅的模方。
因此不同路径先叠加振幅再取模方，交叉项便产生干涉。

可观测量由厄米算符$A$表示，其本征方程为
$A|a\rangle=a|a\rangle$；测量结果是本征值，概率为
$|\langle a|\psi\rangle|^2$。两个算符若不对易，通常不能拥有完备的共同本征态，
这正是不确定关系的结构来源，而不是仪器“不够精密”。

\section{对称性、规范势与微扰}
全同粒子不可由物理手段贴标签。交换两个粒子后，玻色子波函数对称、费米子波函数反对称；
这一简单符号差异决定了凝聚与不相容两种相反的集体行为。

电磁场通过最小耦合$\bm p\to\bm p-q\bm A$进入哈密顿量。规范变换虽可改变$\bm A$
的具体形式，却同时改变波函数相位，闭合回路上的相位差仍可观测。

若$H=H_0+\lambda V$且$\lambda$很小，能量和态可按$\lambda$展开，这称为微扰论。
若多个独立态具有相同未扰能量，称为\textbf{简并}；此时必须先在简并子空间中对角化$V$，
微小外场也可能在一阶就劈裂能级。后文的 Zeeman、Stark 与自旋轨道劈裂都使用这一语言。
\trans{有了“态—算符—对称性—微扰”的统一语言，先从单粒子最基本的量子化、涨落、干涉和隧穿开始。}

\chapter{基础量子效应}
\section{量子化效应}
\Physical
量子化常来自波函数的边界条件：只有能在有限区域内“首尾相容”的驻波才允许存在。

\Model
无限深方势阱$V(x)=0\ (0<x<L)$，阱外$V=\infty$。定态方程为
$-\hbar^2\psi''/(2m)=E\psi$，边界条件$\psi(0)=\psi(L)=0$。

\Derivation
令$k=\sqrt{2mE}/\hbar$，则$\psi=A\sin kx+B\cos kx$。第一边界给$B=0$，
第二边界要求$\sin kL=0$，故$kL=n\pi$，$n=1,2,\ldots$，于是
\begin{equation}E_n=\frac{\hbar^2k_n^2}{2m}
=\frac{n^2\pi^2\hbar^2}{2mL^2},\qquad
\psi_n=\sqrt{\frac2L}\sin\frac{n\pi x}{L}.\end{equation}

对谐振子$H=p^2/(2m)+m\omega^2x^2/2$，定义
\[
a=\sqrt{\frac{m\omega}{2\hbar}}x+\frac{\ii}{\sqrt{2m\hbar\omega}}p,\quad
[a,a^\dagger]=1.
\]
代回得$H=\hbar\omega(a^\dagger a+1/2)$。数算符$N=a^\dagger a$非负，
且$[N,a^\dagger]=a^\dagger$，故从最低态逐级产生$n=0,1,\ldots$：
\begin{equation}E_n=\hbar\omega\left(n+\frac12\right).\end{equation}
氢原子则由库仑势中的径向束缚边界条件得到$E_n\propto-1/n^2$。

\Conclusion
量子化不是“能量天生一格一格”，而是算符谱与物理边界共同筛选允许态。
\Application
量子阱、原子光谱、分子振动；谐振子也是晶格声子和电磁场量子化的基本单元。

\section{零点能效应}
\Physical
基态不能同时具有$x=0$与$p=0$，否则违反非对易性；因而最低能量并非零。
\Derivation
由上节立刻有$E_0=\hbar\omega/2$。也可用
$\Delta x\Delta p\ge\hbar/2$估算。取中心在原点，
\[
E\gtrsim\frac{(\Delta p)^2}{2m}+\frac12m\omega^2(\Delta x)^2
\ge\frac{\hbar^2}{8m(\Delta x)^2}+\frac12m\omega^2(\Delta x)^2.
\]
对$\Delta x$极小化得$\Delta x^2=\hbar/(2m\omega)$，恰有$E_{\min}=\hbar\omega/2$。
\Conclusion
零点能是约束与量子涨落的必然后果；此处不等式被高斯基态饱和。
\Application
分子零点振动、低温比热、真空涨落与 Casimir 效应均与之相连。

\section{不确定性效应}
\Model
令$\delta A=A-\langle A\rangle$、$\delta B=B-\langle B\rangle$，
$\Delta A^2=\langle(\delta A)^2\rangle$。
\Derivation
对$|\alpha\rangle=\delta A|\psi\rangle$与
$|\beta\rangle=\delta B|\psi\rangle$用 Cauchy--Schwarz 不等式：
\[
\Delta A^2\Delta B^2\ge|\langle\alpha|\beta\rangle|^2
\ge [\operatorname{Im}\langle\delta A\delta B\rangle]^2.
\]
因$\operatorname{Im}\langle\delta A\delta B\rangle
=\langle[A,B]\rangle/(2\ii)$，故
\begin{equation}\Delta A\Delta B\ge\frac12|\langle[A,B]\rangle|.\end{equation}
由$[x,p]=\ii\hbar$得到
\begin{equation}\Delta x\Delta p\ge\frac{\hbar}{2}.\end{equation}
\Application
它限制波包压缩、原子稳定尺度和量子测量灵敏度，并直接解释零点能。

\section{叠加与干涉效应}
\Physical
若探测器无法区分粒子来自哪条缝，两条路径不是互斥的经典事件，而是相干振幅。
\Derivation
\begin{align}
\psi&=\psi_1+\psi_2,\\
|\psi|^2&=|\psi_1|^2+|\psi_2|^2+2\operatorname{Re}(\psi_1^*\psi_2).
\end{align}
若$\psi_j=A_j\ee^{\ii\phi_j}$，干涉项为$2A_1A_2\cos(\phi_2-\phi_1)$：
相位差为$2\pi n$时增强，为$(2n+1)\pi$时相消。记录路径信息会使路径与环境纠缠，
约化密度矩阵的非对角元衰减，条纹随之消失。
\Application
电子双缝、中子干涉、原子干涉仪；AB 相位和 Berry 相位最终也通过干涉读出。

\section{隧穿效应}
\Model
矩形势垒$V=V_0$（$0<x<a$），取$E<V_0$，垒内解为
$C\ee^{\kappa x}+D\ee^{-\kappa x}$，其中$\kappa=\sqrt{2m(V_0-E)}/\hbar$。
\Derivation
在$x=0,a$连续匹配$\psi,\psi'$可得
\[
T=\left[1+\frac{V_0^2\sinh^2(\kappa a)}{4E(V_0-E)}\right]^{-1}.
\]
厚势垒$\kappa a\gg1$时
\[
T\simeq\frac{16E(V_0-E)}{V_0^2}\ee^{-2\kappa a}.
\]
对缓变势垒，禁阻区局域波数为
$\kappa(x)=\sqrt{2m[V(x)-E]}/\hbar$。WKB 连接公式给出主指数
\begin{equation}
T\sim\exp\left[-2\int_{x_1}^{x_2}\kappa(x)\,\dd x\right].
\end{equation}
\Conclusion
经典禁阻区中波函数并不立刻为零，而是指数衰减；有限宽度后仍留下透射振幅。
\Application
$\alpha$衰变是核内粒子穿透库仑势垒；STM 电流对针尖距离指数敏感，可成像原子表面。
\trans{单粒子空间波函数之外，电子还带有自旋；把角动量放入磁场，离散投影便转化为可测的束流分裂和谱线劈裂。}

\chapter{自旋、磁场与能级劈裂}
\section{Stern--Gerlach 效应}
\Physical
自旋是内禀角动量，不是经典小球自转。自旋$1/2$在任一轴上只有$\pm\hbar/2$两个投影。
\Model
磁矩与场相互作用
\begin{equation}H=-\bm\mu\cdot\bm B,\qquad \bm\mu=\gamma\bm S.\end{equation}
\Derivation
势能$U=-\bm\mu\cdot\bm B$，故非均匀场中
\begin{equation}F_z=-\partial_zU=\partial_z(\bm\mu\cdot\bm B).\end{equation}
若场梯度选定$z$轴，$S_z$测量值只能为$\pm\hbar/2$，于是力也只有两个主要值，
银原子束分裂成两束，而非经典连续带。
\Application
它奠定自旋测量与量子比特读出的概念基础，并连接 Zeeman 能级分裂。

\section{Zeeman 效应}
\Physical
磁矩在外磁场中的取向能量依赖角动量投影。
\Derivation
忽略自旋时$\bm\mu_L=-e\bm L/(2m)$，对电子且$\bm B=B\hat z$，
\begin{equation}H'=\mu_BB\frac{L_z}{\hbar},\qquad
\Delta E=\mu_BBm_l,\quad \mu_B=\frac{e\hbar}{2m}.\end{equation}
这是正常 Zeeman 效应。加入自旋后，$\bm L,\bm S$先耦合成$\bm J=\bm L+\bm S$，
对固定$J$的角动量投影可得
\begin{equation}\Delta E=\mu_Bg_JBm_J,\quad
g_J=1+\frac{J(J+1)+S(S+1)-L(L+1)}{2J(J+1)}.\end{equation}
不同$g_J$造成更复杂的反常 Zeeman 多重线。
\Application
原子谱线、磁共振与天体磁场测量；强场下过渡到 Paschen--Back 区域。

\section{Paschen--Back 效应}
\Model
取$H_{SO}=A\bm L\cdot\bm S$与
$H_B=\mu_BB(L_z+g_sS_z)/\hbar$。
\Derivation
弱场$\mu_BB\ll |A|\hbar^2$时，$L,S$先耦合，$J,m_J$为好量子数（即相应算符
与主哈密顿量对易、可稳定标记本征态）。强场
\[
\mu_BB\gg |A|\hbar^2
\]
时，外场分别固定$\bm L,\bm S$的投影，$m_l,m_s$成为好量子数，近似能移
\[
\Delta E\simeq\mu_BB(m_l+g_sm_s)+A\hbar^2m_lm_s.
\]
\Conclusion
Paschen--Back 效应是 Zeeman 效应的强场极限，量子数从耦合表象重组为非耦合表象。

\section{自旋轨道耦合效应}
\Physical
电子静止系中，核的电场经相对论变换产生有效磁场；电子自旋磁矩与其耦合。
\Model
\begin{equation}H_{SO}=\xi(r)\bm L\cdot\bm S,\qquad
\xi(r)=\frac{1}{2m^2c^2r}\frac{\dd V}{\dd r},\end{equation}
其中$1/2$包含 Thomas 进动修正。
\Derivation
\begin{equation}\bm L\cdot\bm S=\frac12(J^2-L^2-S^2).\end{equation}
在$|l,s;j,m_j\rangle$中，
\[
\Delta E_{SO}=\frac{\langle\xi\rangle\hbar^2}{2}
[j(j+1)-l(l+1)-s(s+1)],
\]
故同一$n,l$按$j=l\pm1/2$劈裂。
\Application
原子精细结构、半导体自旋电子学与拓扑绝缘体；它也使“自旋 Hall”成为可能。
\trans{磁场耦合角动量；电场则直接耦合电荷分布。随后把视野扩大到光与物质的能量、振动和反冲交换。}

\chapter{电场与光谱效应}
\section{Stark 效应}
\Model
均匀电场取$\bm E=E\hat z$，电荷$q$的微扰为
\begin{equation}H'=-q\bm E\cdot\bm r=-qEz.\end{equation}
\Derivation
非简并微扰给出
\begin{equation}\Delta E_n^{(1)}=\langle n|H'|n\rangle.\end{equation}
有反演对称且宇称确定的态满足$\langle z\rangle=0$，故一阶消失。态的一阶混合代回能量得
\begin{equation}\Delta E_n^{(2)}=\sum_{m\ne n}
\frac{|\langle m|H'|n\rangle|^2}{E_n^{(0)}-E_m^{(0)}}\equiv-\frac12\alpha_nE^2.\end{equation}
氢原子$n\ge2$存在简并，需在简并空间对角化$H'$，因而可出现线性 Stark 劈裂。
\Application
原子极化率、谱线调谐、里德堡原子操控；与 Zeeman 效应同属外场微扰。

\section{Lamb shift（兰姆位移）}
\Physical
Dirac 方程只依赖$n,j$，预言氢原子$2S_{1/2}$与$2P_{1/2}$简并。量子化电磁场的真空涨落
使电子位置抖动，同时电子会发射并重吸收虚光子，产生自能修正；两态在核附近概率不同，
修正遂不同。
\Model
普通非相对论量子力学把电磁场当作给定经典背景，既没有光子真空，也没有电子自能与真空极化，
所以不能从其哈密顿量得到 Lamb shift。QED 经重整化后给出有限的能级差。
\Conclusion
氢的$2S_{1/2}$能级相对$2P_{1/2}$约上移$1057.8\,\mathrm{MHz}$（主导项），
实验是 QED 的经典检验。它与零点涨落、Casimir 效应共享“量子化场真空”这一来源。

\section{Raman 效应}
\Physical
光子经虚中间态与分子振动或晶格声子交换能量，散射光频率因此改变。
\Derivation
能量守恒给
\begin{equation}\hbar\omega_s=\hbar\omega_i\pm\hbar\Omega.\end{equation}
产生一个振动量子时$\omega_s=\omega_i-\Omega$（Stokes）；湮灭已有振动量子时
$\omega_s=\omega_i+\Omega$（anti-Stokes）。后者强度受热布居抑制，二者强度比可测温。
\Application
分子指纹识别、声子谱、应力和温度测量；与弹性 Rayleigh 散射相对。

\section{Franck--Condon 效应}
\Model
Born--Oppenheimer 近似利用核比电子重、运动更慢，将
\[
\Psi(\bm r,\bm R)=\psi_e(\bm r;\bm R)\chi_v(\bm R)
\]
分离。偶极跃迁矩阵元近似分解为电子部分与核振动重叠。
\Derivation
电子跃迁时间内核坐标$\bm R$近似冻结，势能图上的跃迁是“竖直”的。若电子偶极矩对$\bm R$
变化缓慢（Condon 近似），则
\begin{equation}I_{v\to v'}\propto
|\langle\psi'_e|\hat{\bm d}|\psi_e\rangle|^2
|\langle\chi_{v'}|\chi_v\rangle|^2.\end{equation}
第二因子即 Franck--Condon 因子。
\Application
解释分子吸收/荧光振动边带及谱带包络，并连接 Raman 振动光谱。

\section{Mössbauer 效应}
\Physical
自由原子发射动量$p_\gamma=E_\gamma/c$的$\gamma$光子时，反冲能
\begin{equation}E_R=\frac{p_\gamma^2}{2M}\end{equation}
会使发射与吸收线错开。原子束缚在晶体中时，存在不激发声子的零声子过程，动量由整个晶体承担，
其反冲能因总质量巨大而可忽略。
\Model
无反冲概率写成 Lamb--Mössbauer 因子
\[
f=|\langle\text{晶格基态}|\ee^{\ii\bm k\cdot\bm u}|
\text{晶格基态}\rangle|^2=\ee^{-2W},
\]
它与 Debye--Waller 因子同源，温度越低通常越大。
\Application
极窄核共振线用于同质异能移、磁超精细场和引力红移测量。
\trans{前几章主要是单粒子或少体问题。全同粒子的交换对称性一旦与巨量粒子结合，便产生费米海、凝聚与交换关联。}

\chapter{全同粒子、统计与多体效应}
\section{Pauli 不相容效应}
\Derivation
两个费米子的总波函数满足
\begin{equation}\Psi(1,2)=-\Psi(2,1).\end{equation}
若两者处于同一单粒子态$\phi$，交换前后乘积都为$\phi(1)\phi(2)$，于是
$\Psi=-\Psi$，只能有$\Psi=0$。等价地，Slater 行列式出现两列相同而为零。
\Conclusion
两个全同费米子不能占据完全相同的单粒子态；“态”包括空间与自旋全部量子数。
\Application
原子壳层、元素周期表、电子简并压；它是费米简并效应的微观起点。

\section{交换效应}
\Model
占据正交轨道$a,b$的两粒子空间态为
\[
\Psi_\pm=\frac{1}{\sqrt2}[\phi_a(1)\phi_b(2)\pm\phi_b(1)\phi_a(2)].
\]
\Derivation
对对称相互作用$V(1,2)$展开期望值：
\[
\langle V\rangle_\pm=J_{ab}\pm K_{ab},
\]
\[
J_{ab}=\iint |\phi_a(1)|^2V|\phi_b(2)|^2\,\dd1\dd2,\quad
K_{ab}=\iint\phi_a^*(1)\phi_b^*(2)V\phi_b(1)\phi_a(2)\,\dd1\dd2.
\]
$J$是直接库仑项，$K$是交换项；后者来自不可区分的两种指派振幅交叉，并非额外经典力。
\Application
Hartree--Fock、Hund 规则与铁磁交换；它把 Pauli 对称性转化为能量差。

\section{费米简并效应}
\Derivation
体积$V$中的三维自由电子态，每个$\bm k$容纳两个自旋。$T=0$时填满半径$k_F$的费米球：
\[
N=2\frac{V}{(2\pi)^3}\frac{4\pi k_F^3}{3}
=\frac{Vk_F^3}{3\pi^2}.
\]
令$n=N/V$，得$k_F=(3\pi^2n)^{1/3}$及
\begin{equation}E_F=\frac{\hbar^2k_F^2}{2m}
=\frac{\hbar^2}{2m}(3\pi^2n)^{2/3}.\end{equation}
\Conclusion
低温中费米海深处态因 Pauli 阻塞无法跃迁，只有宽约$\kb T$的费米面邻域参与热学与输运。
\Application
金属电子热容、白矮星简并压；与 Kondo 散射中的费米面物理相连。

\section{Bose--Einstein 凝聚效应}
\Model
玻色分布$\bar n(\epsilon)=[\exp(\beta(\epsilon-\mu))-1]^{-1}$，且$\mu\le0$（基态能取零）。
\Derivation
三维自由气体的激发态粒子数上限在$\mu\to0$时为
\[
N_{\rm ex}=V\int_0^\infty
\frac{1}{4\pi^2}\left(\frac{2m}{\hbar^2}\right)^{3/2}
\frac{\sqrt\epsilon\,\dd\epsilon}{\ee^{\beta\epsilon}-1}
=\frac{V}{\lambda_T^3}\zeta(3/2),
\]
其中$\lambda_T=\sqrt{2\pi\hbar^2/(m\kb T)}$。当总粒子数超过此容量，多余粒子必须宏观占据基态。
令$N=N_{\rm ex}(T_c)$得
\begin{equation}T_c=\frac{2\pi\hbar^2}{m\kb}
\left[\frac{n}{\zeta(3/2)}\right]^{2/3}.\end{equation}
\Application
冷原子 BEC、超流；与费米海的“逐态填充”形成鲜明对照。
\trans{周期晶格、磁场、无序和相互作用共同塑造固体中的输运。许多效应不再属于单个电子，而属于整个基态。}

\chapter{凝聚态与量子输运效应}
\section{Hall 效应与整数量子 Hall 效应}
\Physical
经典 Hall 效应来自 Lorentz 力造成横向电荷积累。二维低温强磁场中，轨道运动被量子化，
横向电导形成极精确平台。
\Derivation
取$\bm B=B\hat z$和 Landau 规范$\bm A=(0,Bx,0)$：
\[
H=\frac{1}{2m}[\bm p+e\bm A]^2
=\frac{p_x^2}{2m}+\frac{(p_y+eBx)^2}{2m}.
\]
$p_y$守恒，剩余部分是频率$\omega_c=eB/m$的谐振子，故
\begin{equation}E_n=\hbar\omega_c\left(n+\frac12\right).\end{equation}
每单位面积每个 Landau 能级简并度为$eB/h$，于是填充因子
\begin{equation}\nu=\frac{n_e}{eB/h}=\frac{n_eh}{eB}.\end{equation}
当$\nu$个能级完全填满且费米能落在局域态区间，边缘通道或 Kubo 公式给出
\begin{equation}\sigma_{xy}=\nu\frac{e^2}{h},\qquad \sigma_{xx}\simeq0.\end{equation}
\Application
电阻标准；其拓扑解释是填满能带 Chern 数之和。

\section{分数量子 Hall 效应}
\Model
最低 Landau 能级中动能冻结，电子--电子关联主导。$\nu=1/m$（奇数$m$）的 Laughlin 态为
\begin{equation}
\Psi_m=\prod_{i<j}(z_i-z_j)^m
\exp\left(-\sum_i\frac{|z_i|^2}{4l_B^2}\right),\quad
l_B=\sqrt{\frac{\hbar}{eB}}.
\end{equation}
\Derivation
多项式中每个电子坐标的最高幂约为$m(N-1)$，对应穿过样品的磁通量子数
$N_\Phi\simeq mN$，故$\nu=N/N_\Phi\to1/m$。奇数$m$也保证交换两个电子时波函数反号。
插入一个磁通量子产生准空穴，电荷为$e/m$，并具有任意子统计。
\Application
它是多体拓扑序的范例，与整数效应的单粒子能带拓扑本质不同。

\section{Aharonov--Bohm 效应}
\Model
带电粒子哈密顿量
\[
H=\frac{1}{2m}(\bm p-q\bm A)^2+q\varphi.
\]
\Derivation
沿路径的波函数相对自由情形获得$\exp[\ii q\int\bm A\cdot\dd\bm l/\hbar]$。
两路包围磁通$\Phi$时，由 Stokes 定理
\begin{equation}\Delta\phi=\frac{q}{\hbar}\oint\bm A\cdot\dd\bm l
=\frac{q}{\hbar}\int\bm B\cdot\dd\bm S=\frac{q\Phi}{\hbar}.\end{equation}
即使电子路径处$\bm B=0$，包围区域的磁通仍移动干涉条纹。
\Conclusion
局域可测量的是规范不变的闭路相位；矢势并非纯粹计算辅助量。
\Application
介观环磁阻振荡、超导磁通量子；是 Berry 相位的电磁原型。

\section{Berry 相位效应}
\Physical
绝热是指参数变化慢到系统保持在瞬时非简并本征态中。循环后除动力学相位外，还积累只依赖路径几何的相位。
\Derivation
由
\[
H(\bm R)|n(\bm R)\rangle=E_n(\bm R)|n(\bm R)\rangle
\]
设$|\psi\rangle=\ee^{\ii\alpha_n(t)}|n[\bm R(t)]\rangle$代入 Schrödinger 方程，左乘
$\langle n|$得
\[
\dot\alpha_n=-\frac{E_n}{\hbar}
+\ii\langle n|\nabla_{\bm R}n\rangle\cdot\dot{\bm R}.
\]
闭合回路$C$的几何部分为
\begin{equation}\gamma_n=\ii\oint_C
\langle n(\bm R)|\nabla_{\bm R}n(\bm R)\rangle\cdot\dd\bm R.\end{equation}
\Application
分子几何相位、绝热量子控制；Berry 曲率在 Brillouin 区的积分给 Chern 数和 Hall 电导。

\section{Josephson 效应}
\Model
两个弱耦合超导凝聚态
$\Psi_j=\sqrt{n_j}\ee^{\ii\phi_j}$，满足
$\ii\hbar\dot\Psi_1=\mu_1\Psi_1-K\Psi_2$及交换$1,2$的方程。
\Derivation
把振幅与相位代入并分离实虚部。在$n_1\simeq n_2$时，粒子数流正比于
$\sin(\phi_1-\phi_2)$；乘 Cooper 对电荷$2e$得
\begin{equation}I=I_c\sin\Delta\phi.\end{equation}
电势差使化学势差为$2eV$，从相位方程得到
\begin{equation}\frac{\dd\Delta\phi}{\dd t}=\frac{2eV}{\hbar}.\end{equation}
$V=0$有直流超流；恒定$V$产生频率$2eV/h$的交流电流。
\Application
SQUID、量子电压标准、超导量子比特；本质是宏观相位相干。

\section{Meissner 效应}
\Model
London 方程之一为
\[
\nabla\times\bm j_s=-\frac{n_se^2}{m}\bm B.
\]
\Derivation
结合静磁 Maxwell 方程$\nabla\times\bm B=\mu_0\bm j_s$并用
$\nabla\cdot\bm B=0$：
\[
\nabla\times(\nabla\times\bm B)=-\nabla^2\bm B
=-\frac{\mu_0n_se^2}{m}\bm B,
\]
即
\begin{equation}\nabla^2\bm B=\frac{\bm B}{\lambda_L^2},\qquad
\lambda_L=\sqrt{\frac{m}{\mu_0n_se^2}}.\end{equation}
半无限样品中$B(x)=B(0)\ee^{-x/\lambda_L}$。
\Conclusion
超导体在冷却穿越临界温度时主动排磁；完美导体只冻结初始磁通，二者不同。

\section{Kondo 效应}
\Model
\begin{equation}H=\sum_{k\sigma}\epsilon_kc_{k\sigma}^\dagger c_{k\sigma}
+J\bm S\cdot\bm s(0),\end{equation}
其中反铁磁交换$J>0$耦合杂质自旋与导带电子。
\Derivation
二阶以上的自旋翻转散射对费米面中间态积分，产生
$J_{\rm eff}\sim J+2\rho J^2\ln(D/\kb T)$，故降温时有效耦合增强，电阻出现
$-\ln T$趋势，有限阶微扰最终失效。令耦合变强的尺度定义
\[
\kb T_K\sim D\exp[-1/(2\rho J)].
\]
\Conclusion
$T_K$以下导带电子集体屏蔽杂质自旋，形成多体 Kondo 单态；它不是单次杂质散射所能描述。

\section{Anderson 局域化效应}
\Model
\begin{equation}H=\sum_i\epsilon_ic_i^\dagger c_i
-t\sum_{\langle ij\rangle}(c_i^\dagger c_j+\mathrm{h.c.}),\end{equation}
其中随机$\epsilon_i$描述无序。
\Physical
电子在许多随机路径上多重散射；时间反演路径相干增强返回概率。足够强无序使扩散被干涉抑制，
本征态呈
\begin{equation}|\psi(r)|\sim\ee^{-r/\xi},\end{equation}
$\xi$为局域长度。三维可出现迁移率边和金属--绝缘体转变；一、二维无相互作用体系更易局域。
\Application
无序薄膜、冷原子 speckle 势；与由能隙造成的普通带绝缘不同。

\section{Bloch 振荡与 Wannier--Stark 局域化}
\Derivation
半经典方程
\begin{equation}\hbar\frac{\dd k}{\dd t}=-eE\end{equation}
给$k(t)=k_0-eEt/\hbar$。晶格能带满足$\varepsilon(k+2\pi/a)=\varepsilon(k)$，
群速度$v=(1/\hbar)\partial\varepsilon/\partial k$因而周期变化，周期
$T_B=2\pi\hbar/(eEa)$，即
\begin{equation}\omega_B=\frac{eEa}{\hbar}.\end{equation}
在实空间，邻近晶格位置的势能相差$eEa$，能级形成
$E_n=E_0+neEa$的 Wannier--Stark 梯，理想相干态呈场致局域。
\Application
普通固体中散射太快，半导体超晶格和冷原子晶格中可观测。

\section{Coulomb blockade（库仑阻塞）}
\Model
电容为$C$的小岛增加一个电子的典型充电尺度
\begin{equation}E_C=\frac{e^2}{2C}.\end{equation}
更完整的恒相互作用模型为$E(N)=(-Ne+Q_g)^2/(2C)$。
\Derivation
隧穿若使自由能增加便被抑制。当$\kb T\ll E_C$且结电阻足够大以保持电荷量子化，
低偏压下没有可用的能量完成$N\to N+1$，电流阻塞；门电压调到相邻电荷态简并时出现电导峰。
\Application
单电子晶体管、量子点电荷传感；是能量量子化与电子相互作用共同造成的输运阈值。
\trans{凝聚态把物质波的相干性变成电流；量子光学则直接操控光子与少数能级之间的相干振幅。}

\chapter{量子光学效应}
\section{光电效应}
\Derivation
单个光子能量$h\nu$用于克服逸出功$\phi$，其余成为电子最大动能：
\begin{equation}K_{\max}=h\nu-\phi.\end{equation}
因此截止频率$\nu_0=\phi/h$；遏止电压满足$eV_s=K_{\max}$。
\Conclusion
强度主要改变光电子数，频率决定单电子最大能量，显示光的能量交换具有光子颗粒性。

\section{Compton 效应}
\Derivation
初始电子静止，光子四动量守恒。三维动量给
$\bm p_e=\bm p-\bm p'$，故
$p_e^2=p^2+p'^2-2pp'\cos\theta$。能量守恒
$pc+m_ec^2=p'c+\sqrt{p_e^2c^2+m_e^2c^4}$。
移项平方并代入$p=h/\lambda$，消去$p_e$后得
\begin{equation}\Delta\lambda=\lambda'-\lambda
=\frac{h}{m_ec}(1-\cos\theta).\end{equation}
\Application
证明光子携带动量；用于高能光谱、医学成像背景分析。

\section{受激辐射效应}
\Model
两能级布居$N_1,N_2$与频率$\nu$的辐射能量密度$\rho(\nu)$作用时：
\[
R_{\rm abs}=B_{12}\rho N_1,\quad
R_{\rm sp}=A_{21}N_2,\quad
R_{\rm st}=B_{21}\rho N_2.
\]
热平衡配合 Planck 分布要求$g_1B_{12}=g_2B_{21}$及
$A_{21}/B_{21}=8\pi h\nu^3/c^3$。
\Physical
受激发射的振幅由驱动场锁相，所发光子与入射场同频、同相、同传播模（理想模式描述下）。
若粒子数反转使受激发射超过吸收，谐振腔反馈便产生激光。

\section{Rabi 振荡}
\Model
二能级间隔$\hbar\omega_0$，经典场频率$\omega$。偶极近似下
$H=\hbar\omega_0\sigma_z/2-\bm d\cdot\bm E_0\cos\omega t$。
旋转波近似（丢弃快速反向旋转项）后的旋转框架哈密顿量为
\[
H_{\rm RWA}=\frac{\hbar}{2}
\begin{pmatrix}-\Delta&\Omega_R\\ \Omega_R&\Delta\end{pmatrix},
\quad\Delta=\omega-\omega_0.
\]
\Derivation
其本征值差为$\hbar\Omega$，$\Omega=\sqrt{\Omega_R^2+\Delta^2}$。
由基态出发对矩阵指数化，激发态概率为
\begin{equation}P_e(t)=\frac{\Omega_R^2}{\Omega_R^2+\Delta^2}
\sin^2\left[\frac{\sqrt{\Omega_R^2+\Delta^2}}{2}t\right].\end{equation}
\Application
原子钟、核磁共振、量子门；“二能级系统”即只保留两个相关能态的有效模型。

\section{Autler--Townes 劈裂}
\Physical
强耦合场把裸态$|2,n\rangle$与$|3,n-1\rangle$混合成 dressed states（物质与场的联合本征态）。
共振时有效$2\times2$矩阵的非对角元为$\hbar\Omega_c/2$，本征能量相差
$\hbar\Omega_c$。弱探测光因而看到两个共振峰，频率劈裂尺度约为$\Omega_c$；
失谐时为$\sqrt{\Omega_c^2+\Delta_c^2}$。
\Application
强场谱学；它与 EIT 的谱形可相似，但前者强调真实 dressed 能级分裂，后者强调吸收振幅相消。

\section{EIT（电磁诱导透明）}
\Model
$\Lambda$系统有两个长寿命低态$|1\rangle,|2\rangle$和激发态$|3\rangle$。
探测场以$\Omega_p$耦合$|1\rangle\leftrightarrow|3\rangle$，控制场以$\Omega_c$
耦合$|2\rangle\leftrightarrow|3\rangle$。
\Derivation
双光子共振时，相互作用哈密顿量存在零本征值暗态
\begin{equation}|D\rangle=
\frac{\Omega_c|1\rangle-\Omega_p|2\rangle}
{\sqrt{|\Omega_c|^2+|\Omega_p|^2}}.\end{equation}
它在$|3\rangle$上的振幅严格为零：两条激发路径相消，故无自发辐射并对探测光透明。
\Application
慢光、量子存储；核心是概率振幅干涉，而非简单饱和吸收。

\section{Hanbury Brown--Twiss 效应与反聚束}
\Model
量子二阶关联函数为
\begin{equation}g^{(2)}(\tau)=
\frac{\langle:\hat I(t)\hat I(t+\tau):\rangle}{\langle\hat I\rangle^2}
=\frac{\langle a^\dagger(t)a^\dagger(t+\tau)a(t+\tau)a(t)\rangle}
{\langle a^\dagger a\rangle^2}.\end{equation}
\Conclusion
热光$g^{(2)}(0)=2$（聚束），理想相干光$=1$（Poisson 统计），理想单光子态$=0$。
$g^{(2)}(0)<1$表示一次探测后短时间内再次探测的概率被压低，称反聚束；经典正定强度随机过程
不能产生这一特征。
\Application
恒星角直径测量、单光子源认证和量子发光体动力学。
\trans{前述理论大多固定粒子数或把光场作近似。相对论与量子场论允许正负能量、真空极化乃至粒子产生。}

\chapter{相对论量子力学与量子场论相关效应}
\section{Klein 隧穿效应}
\Physical
Dirac 方程的一阶空间结构和正负能谱使高势垒内部可连接到“空穴”分支，故势垒越高并不总意味着
指数衰减。场论语言中，这与粒子--反粒子通道有关。
\Application
石墨烯低能载流子满足二维无质量 Dirac 哈密顿量
$H=v_F\bm\sigma\cdot\bm p+V(x)$。平滑势垒正入射时赝自旋（两个子晶格分量）
守恒，反射要求赝自旋翻转而被禁止，理想情形$T(0)=1$。斜入射或开隙会降低透射。

\section{Zitterbewegung（颤动效应）}
\Physical
Dirac 速度算符为$c\bm\alpha$，不与自由哈密顿量对易。Heisenberg 图像中的位置包含一个
频率约$2E/\hbar$、幅度约 Compton 波长的快速振荡项。它来自波包中正能与负能分量的相对相位演化；
若严格投影到单一能量分支，交叉干涉项消失。
\Application
基本电子的直接观测受限于极短尺度，但离子阱、冷原子和固体 Dirac 材料可模拟。

\section{Casimir 效应}
\Physical
每个电磁模具有$\hbar\omega/2$零点能。两平行理想导体板限制法向允许波数，改变板内外模密度。
将两种构型的发散零点能作正则化差值，有限且依赖板距$a$。
\Conclusion
零温理想板的单位面积能量$E/A=-\pi^2\hbar c/(720a^3)$，故
\begin{equation}\frac{F}{A}=-\frac{\partial(E/A)}{\partial a}
=-\frac{\pi^2\hbar c}{240a^4}.\end{equation}
负号表示吸引。真实材料需加入有限电导、粗糙度与温度修正。

\section{Schwinger 效应}
\Physical
强电场把真空中的正负能谱倾斜，使虚粒子对可穿越质量能隙势垒而成为实粒子。世界线或 WKB
作用量给出非微扰指数
\begin{equation}\Gamma\propto
\exp\left(-\frac{\pi m^2c^3}{qE\hbar}\right).\end{equation}
临界尺度$E_S=m^2c^3/(q\hbar)$极大；指数关于$1/E$非解析，因此任何有限阶通常微扰都得不到它。
\Application
强场 QED、超强激光；与 Klein 隧穿共享 Dirac 谱中的跨能隙隧穿图像。
\trans{最后，Berry 几何在整个动量空间中积累成整数或$\mathbb Z_2$不变量，使边界导电不依赖微观细节。}

\chapter{拓扑量子效应}
\section{量子自旋 Hall 效应}
\Physical
自旋轨道耦合可使相反自旋获得相反的有效横向响应，形成一对反向传播的 helical edge states：
运动方向与自旋锁定。系统总体保持时间反演对称，不需要外磁场，净电荷 Hall 电导通常为零。
\Conclusion
普通量子 Hall 边缘态是同向传播（手性）并破坏时间反演；量子自旋 Hall 的 Kramers 对在保持
时间反演的非磁无序下不能彼此弹性背散射，由$\mathbb Z_2$拓扑不变量保护。

\section{反常量子 Hall 效应}
\Model
磁性与自旋轨道耦合可在零外磁场下产生非零 Berry 曲率$\Omega_n(\bm k)$。填满二维能带的
\[
C=\frac{1}{2\pi}\sum_{n\in{\rm occ}}\int_{\rm BZ}\Omega_n(\bm k)\,\dd^2k
\]
为整数 Chern 数。
\Conclusion
\begin{equation}\sigma_{xy}=C\frac{e^2}{h}.\end{equation}
量子化来自整个 Brillouin 区波函数的拓扑绕卷，而非净外磁通；磁序负责破坏时间反演。
\Application
零场手性边缘输运、低功耗拓扑器件研究。

\section{拓扑绝缘体边缘态效应}
\Physical
拓扑绝缘体体内有能隙，边界却导电。把拓扑非平凡材料与真空（平凡绝缘体）相接，拓扑不变量必须
改变；只要保护对称性与体能隙不被破坏，这个改变要求边界能隙闭合，产生边界态。
\Conclusion
这就是体--边对应。Chern 绝缘体中，$C$之差等于净手性边缘通道数；二维/三维时间反演拓扑绝缘体
由$\mathbb Z_2$不变量区分，边缘/表面出现奇数对 Kramers/Dirac 态。Berry 联络与曲率是计算这些
不变量的语言，拓扑保护意味着弱无序可改变色散，却不能单独移除全部边界态。
\Application
拓扑表面输运、自旋电子学，以及与超导邻近效应结合的拓扑量子计算探索。

\chapter{总结：纷繁效应背后的统一结构}
\begin{longtable}{p{0.28\textwidth}p{0.64\textwidth}}
\toprule
统一机制 & 代表效应与逻辑\\
\midrule
\endhead
边界条件导致量子化 & 方势阱、谐振子与 Landau 能级都是允许本征模的离散谱。\\
对易关系导致不确定性 & $[x,p]=\ii\hbar$使相空间不能同时无限压缩，并留下零点能。\\
概率振幅叠加导致干涉 & 双缝、EIT、AB 相位都由不可区分路径的交叉项显现。\\
禁阻区指数衰减导致隧穿 & 普通 WKB、$\alpha$衰变、Klein 与 Schwinger 效应是不同能谱中的越垒。\\
角动量与外场耦合 & Stern--Gerlach、Zeeman、Paschen--Back 与自旋轨道耦合产生投影选择和劈裂。\\
交换对称性决定统计 & 费米反对称性产生 Pauli 原理和费米海；玻色对称性允许宏观凝聚。\\
规范势改变相位 & AB 效应表明无局域场的路径仍可由磁通获得规范不变相位。\\
绝热演化产生几何相位 & Berry 相位的局域曲率在参数空间积分，进一步形成拓扑不变量。\\
多体关联创造新准粒子 & Kondo 屏蔽和分数量子 Hall 态不能还原为独立电子，后者产生分数电荷。\\
拓扑不变量保护边界态 & Chern 数或$\mathbb Z_2$不变量通过体--边对应约束稳健边缘输运。\\
\bottomrule
\end{longtable}

这些结构并非彼此隔绝。例如 Landau 量子化先提供离散能级，交换统计决定填充方式，
无序帮助形成平台，多体关联产生分数态，Berry 曲率最终解释量子化电导。学习一个新效应时，
最有效的问题通常不是“公式是什么”，而是：相关自由度是什么？哪些振幅相加？什么对称性被保持或破坏？
谱隙由何而来？实验实际测量的是能量、相位、关联还是输运？

\chapter*{推荐阅读}
\addcontentsline{toc}{chapter}{推荐阅读}
\begin{enumerate}
\item D. J. Griffiths and D. F. Schroeter, \emph{Introduction to Quantum Mechanics}.
\item J. J. Sakurai and J. Napolitano, \emph{Modern Quantum Mechanics}.
\item R. Shankar, \emph{Principles of Quantum Mechanics}.
\item C. Cohen-Tannoudji, B. Diu and F. Laloë, \emph{Quantum Mechanics}.
\item N. W. Ashcroft and N. D. Mermin, \emph{Solid State Physics}.
\item G. D. Mahan, \emph{Many-Particle Physics}.
\item M. O. Scully and M. S. Zubairy, \emph{Quantum Optics}.
\item X.-G. Wen, \emph{Quantum Field Theory of Many-Body Systems}.
\end{enumerate}

\end{document}
